\section{Implementation Details}
\label{app:impl}

\subsection{Network architectures}

\paragraph{Policy}
The policy network (\texttt{BusEmbeddingPolicy}) consists of a shared embedding layer $\phi$ over five categorical descriptors (\texttt{line\_id}, \texttt{bus\_id}, \texttt{station\_id}, \texttt{time\_period}, \texttt{direction}) each projected to a $48$-dim embedding, concatenated with $12$ continuous features, and passed through a two-layer MLP with $256$ hidden units and ReLU activation.
The output maps to a two-dimensional Gaussian action with mean $\mu \in [-1, 1]^2$ (tanh-squashed) and state-independent log-std.
An \texttt{action\_range} scalar of $1.0$ rescales outputs to the operational range; layer normalisation and dropout $0.05$ are applied inside the embedding.

\paragraph{Q ensemble}
Five critic heads share the embedding architecture with independent MLPs of the same width.
Heads are trained with independent minibatches and independent target networks, updated by Polyak averaging with $\tau = 0.005$.

\paragraph{Contrastive discriminator}
The score $f_\psi(s, s')$ is a two-layer MLP over the concatenated $(s, s')$ vector with hidden width $128$, ReLU activation, and a scalar output.
InfoNCE temperature is fixed to $0.1$; $K = \text{batch\_size} - 1$ negatives are drawn in-batch from the opposite-domain samples.

\subsection{Hyperparameters}

\begin{table*}[t]
\centering
\small
\resizebox{\textwidth}{!}{%
\begin{tabular}{ll}
\toprule
Parameter & Value \\
\midrule
Batch size                           & 2048 \\
Offline/online mix ratio             & 1.5 (offline) : 1.0 (online) \\
Ensemble size $\Ecnt$                & 5 \\
LCB coefficient $\beta_{\text{LCB}}$ & 1.0 \\
Offline-Q floor (scalar per minibatch)  & $\E_{\mathcal{B}_{\mathrm{off}}}[\mu_Q(s, a)]$ \\
RE-SAC indep/min blend ratio $\kappa$ & 0.8 \\
Q-std clip coefficient $c$           & 0.5 \\
LCB normalisation                    & by $\max(|\mu_Q|, 1)$ \\
EMA policy decay $\tau$              & $5 \times 10^{-3}$ \\
Policy anchor weight $\lambda_{\text{anc}}$ & 0.1 \\
IS weight clip $[\underline w, \overline w]$ & $[0.1, 10]$ \\
IS hybrid weight $\lambda$           & 1.0 \\
KL anchor coefficient                & 0.5 (when enabled) \\
Offline-only warmup epochs $N_{\text{warm}}$ & 20 (when enabled; distinct from discriminator warmup of 5\,000 grad steps) \\
Updates per epoch                    & 100 (SAC) + 100 (rollout events) \\
Number of epochs                     & 200 \\
Discount $\gamma$                    & 0.99 \\
Policy learning rate                 & $3\mathrm{e}{-4}$ \\
Q learning rate                      & $3\mathrm{e}{-4}$ \\
Discriminator learning rate          & $3\mathrm{e}{-4}$ \\
Polyak $\tau$                        & 0.005 \\
Seeds                                & $\{42, 123, 456, 789, 2024\}$ \\
\bottomrule
\end{tabular}%
}
\caption{Hyperparameters used across all configurations unless explicitly ablated.}
\label{tab:hparams}
\end{table*}

\subsection{Transit-specific data pipeline fixes}
\label{app:impl:pipeline}

Three implementation issues---individually minor but jointly determining whether sim and SUMO data are comparable at all---must be handled before any H2O method can be meaningfully trained on transit data.
We document them here for reproducibility and as a checklist for future applications of H2O methods to new transit benchmarks.

\paragraph{(i) Trajectory-dict isolation}
In the original data collection pipeline, RL-controlled hold time was appended to the same per-tick trajectory buffer used for headway measurement, leaking hold duration into the next bus's \texttt{headway\_prev} feature.
The fix is to keep hold-time annotations in a separate buffer consumed only by the logger, not by the observation constructor.
Without this fix, the effective state distribution shifts between episodes with and without RL intervention, invalidating the Bellman equation.

\paragraph{(ii) Segment-speed rescaling}
The sim-core travel-time model internally reports speeds at the base simulation step (0.5\,s), while SUMO is configured with a doubled step length (1.0\,s).
Multiplying sim-core segment speeds by $2.0$ aligns the two under a shared $\Delta t$.
Without this fix, sim headway predictions differ from SUMO by a $13.2\sigma$ gap on matched initial conditions---large enough to dominate any learned dynamics correction.

\paragraph{(iii) Observation-normaliser freezing}
The data-collection workers called the observation normaliser with default \texttt{update=True}, overwriting the running statistics with collected-distribution statistics rather than preserving those from the offline distribution.
The fix is to pass \texttt{update=False} in collection workers and update statistics only in the offline pre-training loop.
Without this fix, the mean per-step reward on SAC-collected data drifts from $-56$ to $-197$---an order-of-magnitude artefact that masks all algorithmic effects.

\subsection{Offline dataset details}

$\Doff$ aggregates transitions from four behaviour policies, collected under matched SUMO seeds:
(a) SAC-trained: $180{,}000$ transitions from a previously trained online SAC agent, providing near-optimal coverage;
(b) Heuristic headway-threshold: $160{,}000$ transitions from an analytical rule-based controller;
(c) Weak $\epsilon$-random: $180{,}000$ transitions from an $\epsilon$-greedy random policy with $\epsilon = 0.3$, providing coverage of low-quality actions;
(d) RE-SAC SUMO expert ep39: $155{,}000$ transitions from the best checkpoint of a prior RE-SAC run in the SUMO environment (a deterministic learned policy, not a hand-coded rule).
Observations are z-score normalised using full-$\Doff$ statistics; actions are already in $[-1, 1]$ by the policy mapping; rewards are used raw (not normalised).

\subsection{Auto-parallel launcher}
\label{app:impl:parallel}

The experiment sweep uses a hardware-aware launcher \texttt{run\_\allowbreak full\_\allowbreak experiments.sh --parallel auto} that:
(i) detects CPU cores, system RAM, and GPU count/VRAM;
(ii) assumes per-job resource consumption of 2 CPU cores, 4\,GB RAM, and 2\,GB VRAM;
(iii) caps parallelism at 2 jobs per GPU to avoid thrashing, and at a global ceiling of 8;
(iv) forces the SUMO-online group to run sequentially due to libsumo's single-session restriction;
(v) assigns jobs to GPUs in round-robin when multiple GPUs are available.
On our evaluation hardware (32 CPU cores, 29\,GB RAM, $1\times$ 7\,GB GPU), the launcher runs 2 jobs in parallel for SIM-online groups and 1 job at a time for SUMO-online, yielding a total wall time of approximately 120 CPU-hours for the full sweep ($\approx$30 configurations $\times$ 4--5 seeds $\approx$ 150 runs).

\subsection{Action-invariance test details}
\label{app:assumption1}

The data-level test of \Cref{ass:action_invariance} reported in \Cref{sec:exp:analysis} is implemented in \texttt{SimpleSAC/verify\_assumption1.py}.
We summarise its protocol and limitations here.

\paragraph{Inputs and grid}
The real-side h5 is the merged offline buffer ($n_{\text{real}} = 50{,}000$ sub-sampled transitions from $\Doff$).
The sim-side h5 is generated by replaying $\Doff$'s initial-state snapshots through sim-core under uniform-random actions ($n_{\text{sim}} = 12{,}999$ transitions; smaller because some snapshots fail to advance under random actions).
We z-score-normalise the concatenated $(s, s')$ vectors, project to two PCA components, bin into a $5 \times 5$ grid, and discretise actions into $K_a = 3$ bins (no-hold, short-hold $> 0$\,s, long-hold $\geq 30$\,s).

\paragraph{Coverage caveat}
Of $25$ $(s, s')$ cells $\times$ $3$ action bins $= 75$ cell-action combinations, only $6$ have $\geq 20$ samples in both real and sim h5; only $2$ of the $25$ $(s, s')$ cells have all three action bins covered.
The reported $\hat{r} = 0.156$ is therefore computed on a small valid subset, and a bootstrap CI on the same subset is degenerate (all resamples produce the same point estimate because there is essentially one informative pair of cells).
The $\hat{r} = 0.169$ value at the $6 \times 6$ grid (3 of 36 cells valid) is consistent.
We read both numbers as a directionally supportive signal --- variation in $\log \hat\rho$ across actions is small relative to variation across $(s, s')$ in the cells we can measure --- not as a tight statistical verification of \Cref{ass:action_invariance} on the entire state space.
A camera-ready version should rerun the test with (i) larger $n_{\text{sim}}$ ($\geq 50$K, by extending the rollout-for-verification budget), (ii) a coarser PCA grid ($3 \times 3$) for more cells with valid action coverage, and (iii) explicit per-cell sample-count reporting alongside the ratio.

\paragraph{Heatmap}
\Cref{fig:assumption1_heatmap} visualises $\log \hat\rho$ per cell, broken out by action bin.
NaN cells (white) lack sufficient samples; coloured cells span $[-2, 2]$ in $\log \hat\rho$.
The two valid $(s, s')$ cells show consistent log-ratios across all three action bins, which is the qualitative pattern \Cref{ass:action_invariance} predicts.

\begin{figure}[h]
\centering
\includegraphics[width=0.95\linewidth]{figures/assumption1_heatmap.pdf}
\caption{Per-cell $\log \hat{\rho} = \log[\hat{P}_{\text{real}}(s' \mid s, a) / \hat{P}_{\text{sim}}(s' \mid s, a)]$ across a $5 \times 5$ PCA grid in $(s, s')$, broken out by action bin (no-hold / short-hold / long-hold). White cells lack the $\geq 20$-sample threshold needed for a stable estimate; coverage is sparse, and the action-dependence ratio reported in \Cref{sec:exp:analysis} is computed from the small number of cells with valid data in all action bins.}
\label{fig:assumption1_heatmap}
\end{figure}
